\documentclass[11pt,a4paper,sans,english,colorlinks,dvipsnames]{moderncv}        % possible options include font size ('10pt', '11pt' and '12pt'), paper size ('a4paper', 'letterpaper', 'a5paper', 'legalpaper', 'executivepaper' and 'landscape') and font family ('sans' and 'roman')

\moderncvstyle{classic}  
\moderncvicons{awesome}
% style options are 'casual' (default), 'classic', 'oldstyle' and 'banking'
\moderncvcolor{red}                               % color options 'blue' (default), 'orange', 'green', 'red', 'purple', 'grey' and 'black'
%\nopagenumbers{}                                  % uncomment to suppress automatic page numbering for CVs longer than one page
\usepackage[utf8]{inputenc}                       % if you are not using xelatex ou lualatex, replace by the encoding you are using
\usepackage[top=1in,bottom=1in,scale=0.75,a4paper]{geometry}
\usepackage{babel}
\usepackage[bottom]{footmisc}
\usepackage{eurosym}
\definecolor{blue1}{RGB}{56,115,178}
%\renewcommand*{\orcidsocialsymbol}{{\aiSkype~}}
%----------------------------------------------------------------------------------
%            personal data
%----------------------------------------------------------------------------------
%\vspace{-10mm}
\firstname{Shyam}
\familyname{Harimohan Menon}
%\title{Resumé}                               % optional, remove/comment the line if not wanted
\address{Center for Computational Astrophysics}{Flatiron Institute}{NY 10010, USA}         % optional, remove/comment the line if not wanted; the "country" arguments can be omitted or provided empty
\email{smenon@flatironinstitute.org}                               % optional, remove/comment the line if not wanted
\social[github]{shm-1996}
\social[orcid]{0000-0001-5944-291X}
\homepage{shm-1996.github.io}                         % optional, remove/comment the line if not wanted
%\extrainfo{additional information}                 % optional, remove/comment the line if not wanted
% \photo[80pt][0pt]{Suit}                       % optional, uncomment the line if wanted; '64pt' is the height the picture must be resized to, 0.4pt is the thickness of the frame around it (put it to 0pt for no frame) and 'picture' is the name of the picture file
%\quote{Final year Masters student in the Department of Chemistry, Indian Institute of Technology (BHU), Varanasi}                                 % optional, remove/comment the line if not wanted
%
\begin{document}
\hypersetup{
urlcolor=blue}
\urlstyle{same}
%-----       resume       ---------------------------------------------------------
%\vspace{-10mm}
\makecvtitle
\vspace{-10mm}
%\hline
\section{Professional Experience}
% \cventry{Sep 2023 -- present}{Joint Rutgers-Flatiron Postdoctoral Fellow}{}{\newline Center for Computational Astrophysics (CCA), Flatiron Institute, NY, USA \& Dept.\ of Physics and Astronomy, Rutgers University, NJ, USA}{}{}
\cventry{Sep 2025 -- present}{Flatiron Research Fellow (Joint Rutgers--Flatiron Fellow)}{}{\newline Center for Computational Astrophysics (CCA), Flatiron Institute, NY, USA}{}{}
\cventry{Sep 2023 -- Aug 2025}{Postdoctoral Fellow (Joint Rutgers--Flatiron Fellow)}{}{\newline Dept.\ of Physics and Astronomy, Rutgers University, NJ, USA}{}{}

\section{Education}
\cventry{Oct 2019 -- Sep 2023}{Doctor of Philosophy}{Astronomy \& Astrophysics}{\newline Research School of Astronomy \& Astrophysics, Australian National University}{Canberra, Australia}{Thesis: \href{http://hdl.handle.net/1885/299505}{The interplay of star formation and radiation feedback in the ISM}. \newline \textit{Awarded the \href{https://asa.astronomy.org.au/prizes_and-grants/prizes-awards/charlene-heisler-prize/}{Charlene Heisler Prize}}. \newline Advisors: Prof.~C.~Federrath \& Prof.~M.~Krumholz}
\cventry{Aug 2014 -- June 2019}{Integrated Bachelor and Master of Technology (Hons.)}{Engineering Physics}{\newline Indian Institute of Technology (BHU)}{Varanasi, India}{Thesis: \href{https://www.dropbox.com/sh/afarzv6v6pvqyes/AACdRiKa9tuUTN-Z_GbyU1Qfa?dl=0}{HI Fragmentation in the Turbulent Interstellar Medium}.\newline Advisor: A/Prof.~P.~Dutta}

% \section{Research Projects}
% \cventry{August 2020 -- present}{Impact of Radiation Feedback on Star Cluster Formation}{}{\newline Primary Collaborators:A/Prof.~Christoph Federrath \& Prof.~Mark Krumholz }{}{$\bullet$ Developed \texttt{VETTAM}\footnotemark, a novel radiation-hydrodynamics module with Adaptive-Mesh Refinement in the \texttt{FLASH} code, to study the effects of radiation on star cluster formation. [Publication \#~\ref{pubvet}] \\ $\bullet$ Using \texttt{VETTAM}, we tested the role of radiation pressure from dust-reprocessed infrared photons on regulating star formation and driving winds in super-star cluster environments. [Publication \#~\ref{pubIR}]  \\ $\bullet$
% Currently testing the cumulative effects of radiative feedback mechanisms using multiband RHD simulations, and constraining the conditions in which feedback fails to regulate star formation.}
% \footnotetext{Variable Eddington Tensor closed Transport on Adaptive Meshes (VETTAM)}
% \cventry{August 2019 -- May 2021}{Impact of ionising radiation feedback on the turbulent interstellar medium (ISM)}{}{\newline Primary Collaborators: A/Prof.~Christoph~Federrath, Dr.~Rolf Kuiper, Dr.~Pamela Klaassen, Piyush Sharda}{}{$\bullet$ Probed the effects of ionising radiation on driving compressive turbulence and potentially triggering star formation in the ISM using numerical simulations. [Publication \#~\ref{pubpluto}] 
% \\ $\bullet$ Tested the hypothesis of the above study observationally by quantifying the bulk properties of the turbulence in the Carina nebula. [Publication \#~\ref{pubcarina}] \\ $\bullet$ Extended the methodology of the previous study to quantify the dynamical effects of physical mechanisms on turbulence in the LMC. [Publication \#~\ref{publmc}] \\ $\bullet$ Currently studying the role of ionising radiation in shaping the magnetic fields in star-forming gas.}

% \cventry{September 2020 -- Present}{The Spatial Distribution of Star Clusters in External Galaxies}{}{\newline Primary Collaborators: Dr.~Kathryn Grasha, Prof.~Bruce G.~Elmegreen, Prof.~Daniela Calzetti, Dr.~Angela Adamo}{}{$\bullet$ Investigated the fractal distribution of young star clusters in 12 local galaxies with the LEGUS survey. [Publication \#~\ref{publegus}] \\ $\bullet$ Currently extending work to probe the star cluster distributions in local high-redshift analogue galaxies with the HiPEEC survey, and a larger sample of local galaxies with the PHANGS survey.}

\section{Honors \& Grants}
\cventry{July 2024}{Charlene Heisler Prize}{\newline Awarded by the Astronomical Society of Australia for the most outstanding PhD thesis accepted by an Australian university}{}{}{}
\cventry{Feb 2023}{Joint 4-year Rutgers-Flatiron Postdoctoral Fellowship ($\sim$\$400,000 {})}{}{}{}{}
\cventry{Feb 2023}{CCAPP Postdoctoral Fellowship (declined)}{}{}{}{}
\cventry{Feb 2023}{CITA National Fellowship (declined)}{}{}{}{}
\cventry{March 2022}{RSAA Travel Grant ($\sim$\$6000{})}{}{}{}{}
\cventry{Oct 2021}{ANU Olin J Eggen Research Award ($\sim$\$2000{})}{\newline Awarded to a PhD student for excellence in research}{}{}{}
\cventry{Oct 2019}{ANU International PhD Scholarship ($\sim$\$80,000{})}{}{}{}{}
\cventry{Oct 2019}{ANU RSAA Supplementary Research Scholarship ($\sim$\$7000{})}{}{}{}{}
\cventry{May 2017}{Mitacs Globalink Research Internship Award ($\sim$\$5000{})}{}{}{}{}
\cventry{April 2013}{Kishore Vaigyanik Protsahan Yojana (KVPY) Scholarship}{}{}{}{}

\newpage
\section{Leadership and Scientific Service}

\cventry{2026-2028}{Leader: International Space Science Institute (ISSI) International Team}{\newline Leading the competitively-selected ISSI International Team \#26-690 focused on multiscale modeling of dust in the interstellar medium of galaxies in the early Universe ($z\gtrsim 6$).}{}{}{}
\cventry{Feb 2027}{Lead Organiser, Scientific Workshop on "The Lifecycle of Dense Star Clusters"}{\newline Lead organiser of a focused workshop consisting of $\sim$50 international experts on dense star cluster formation and evolution in New York City}{}{}{}
\cventry{April 2026}{Expert grant panel reviewer: National Science Foundation (NSF) Grants}{\newline Invited to provide expert reviews on an NSF-AAG grant panel.}{}{}{}
\cventry{July 2021 -- present}{Reviewer: ApJ, ApJL, MNRAS, J. Comput. Phys., JOAA, Open Journal of Astrophysics}{}{}{}{}
\cventry{Feb 2022 -- May 2023}{Chair, Seminars Committee, RSAA}{\newline Heading the seminars team; responsible for organising, hosting and inviting speakers for colloquia}{}{}{}
\cventry{Sept 2020 -- Feb 2022}{Member of Seminars Committee, RSAA}{\newline Hosted >30 external colloquia and seminars at RSAA}{}{}{}

\section{Students Mentored}
\textbf{\textit{2 PhD students + 2 master's + 2 undergraduate students}} \newline

\cventry{Jan 2025 -- May 2025}{Thummim Mekuria, CCA Predoctoral Project (PhD student at UC Santa Cruz)}{\newline Project title: The Feedback-Matter Interplay in High-z  Compact Starbursts}{}{}{}

\cventry{Jan 2025 -- present}{Avi Chen, PhD student, Rutgers University}{\newline Project title: \href{https://iopscience.iop.org/article/10.3847/1538-4357/ae969f}{Role of Radiative Feedback in Population III Star Formation}}{}{}{}

\cventry{Nov 2024 -- Nov 2025}{Lisanne van Veenen, Master's student at Leiden University}{\newline Project title: \href{https://www.aanda.org/articles/aa/full_html/2026/04/aa58114-25/aa58114-25.html}{Radiation magnetohydrodynamics simulations of Population III star formation during the epoch of reionization}}{}{}{}

\cventry{March 2024 -- March 2025}{Dhavala Sai Srinivas, Master's student at IISER-Kolkata}{\newline Project title: The escape fraction of ionising photons in magnetized turbulent media}{}{}{}

\cventry{Jan 2025-- March 2025}{Frances Mathias, Undergraduate, Rutgers University}{\newline Project title: Stochasticity in Stellar Population Synthesis}{}{}{}

\cventry{Jan 2022 -- March 2022}{Hank Hua, undergraduate student at ANU (co-supervised)}{\newline Project title: Effects of spiral arms on the Two Point Correlation Function}{}{}{}


\section{Teaching Experience}

\cventry{Oct 2026}{Guest Lecture, Physics 345 -- Introduction to Computational Astrophysics}{\newline Guest lecture on high performance computing for graduate students}{}{}{}

\cventry{Sep 2025}{Guest Lecture, Physics 610 -- Interstellar Matter}{\newline Guest lectures on ISM physics for graduate students}{}{}{}

\cventry{Nov 2024}{Guest Lecture, Physics 345 -- Introduction to Computational Astrophysics}{\newline Guest lecture on parallel computing for graduate students}{}{}{}


\cventry{Oct 2022}{Guest Lecture, ASTR 4012 -- Astrophysical Gas Dynamics}{\newline Guest Lecture on The Physics of Radiation Hydrodynamics for graduate students}{}{}{}

\cventry{Feb 2022 -- May 2023}{Teaching Assistant, ASTR 1003 -- Astronomy \& Space}{\newline Respond to undergraduate student queries on the course forum and assist in tutorials}{}{}{}

\section{Outreach Activities}
\cventry{Sep 2026}{Public talk, \emph{Mortal Matters, Greenwood Public Cemetery}}{\newline Delivered an invited public talk on "All Ends are Beginnings" at the \href{https://www.green-wood.com/event/mortal-matters/}{"Mortal Matters"} event held at the Greenwood Public Cemetery.}{}{}{}
\cventry{Aug 2026}{Public talk, \emph{Astronomy on Tap (AoT), Bengaluru}}{\newline Delivered an invited AoT public outreach talk on "Survival of the Fittest: Lessons from Ancient Stellar Nurseries" that obtained \href{https://www.instagram.com/p/Db-BSjjNAPz/?stkn=MTNxMWVmd3hsbnhxdg==}{press coverage}.}{}{}{}
\cventry{Sep 2024 -- May 2025}{STEM professional volunteer: Letters to a Pre-Scientist}{\newline Program to exchange letters with middle school students to broaden exposure to the day-to-day of STEM researchers.}{}{}{}
\cventry{July 2021 -- Dec 2022}{Hosting public nights, Mount Stromlo Observatory}{\newline Hosting telescope viewing nights at the observatory for public school students and adults as part of the outreach team at ANU.}{}{}{}

% \section{Proposals}
% \cventry{April 2026}{Co-Investigator, \normalfont JWST-GO-10202
% : MIRAGE: MIRI 
% Investigation of Red and Ancient Galaxies in the Early Universe}{}{}{}{}
% \cventry{July 2025}{Co-Investigator, \normalfont ALMA Cycle 12 Proposal 2025.1.00874.S: Feeding and Feedback: Connecting Molecular Gas to the Young Stellar Clusters of M83}{}{}{}{}

\section{Invited and Prize Talks}

\cventry{April 2027}{Deepdish-Baguette: Stars, Black Holes, and Feedback From Cosmic Dawn to Cosmic Noon}{}{}{}{}
\cventry{Nov 2026}{Department of Physics Colloquium: University of Texas, Dallas}{}{}{}{}
\cventry{Sep 2026}{Colloquium: Tata Institute of Fundamental Research, Mumbai}{}{}{}{}
\cventry{Sep 2026}{Seminar: Raman Research Institute (RRI), Bengaluru}{}{}{}{}
\cventry{May 2026}{Astrophysics Seminar: American Museum of Natural History}{}{}{}{}
\cventry{April 2026}{Star Formation/ISM Rendezvous: Princeton University}{}{}{}{}
\cventry{Jan 2026}{The growth of galaxies in the early Universe XI, Sesto, Italy}{}{}{}{}
\cventry{Dec 2025}{Theoretical Astrophysics Seminar, UC Berkeley}{}{}{}{}
\cventry{Aug 2025}{Santa Cruz Galaxy Workshop, UC Santa Cruz}{}{}{}{}
\cventry{July 2025}{Review Talk, Numerical Recipes in Star Formation, Lorentz Center, Leiden University}{}{}{}{}
\cventry{Aug 2024}{Cosmic Dawn Revealed by JWST, Kavli Institute of Theoretical Physics}{}{}{}{}
\cventry{June 2024}{Charlene Heisler Prize Talk, ASA Annual Meeting}{}{}{}{}
\cventry{March 2024}{ISM Salon, CCA Flatiron Institute}{}{}{}{}
\cventry{March 2024}{Radiation Transport in Astrophysics, Heidelberg}{}{}{}{}
\cventry{July 2023}{Seminar: Inter-University Centre for Astronomy and Astrophysics (IUCAA), Pune}{}{}{}{}
\cventry{June 2022}{Seminar: Centre de Recherche Astrophysique de Lyon (CRAL)}{}{}{}{}
\cventry{Dec. 2021}{Colloquium: Indian Institute of Astrophysics, Bangalore}{}{}{}{}
\cventry{March 2020}{Modelling High-Mass Stellar Feedback Workshop, T\"{u}bingen, Germany}{}{}{}{}


\section{Contributed Talks (\emph{selected})}
\cventry{July 2026}{Cake Talk: Cosmic Dawn Center, Copenhagen}{}{}{}{}
\cventry{June 2026}{EAS Annual Meeting (Special Session 5 and 7), Lausanne, Switzerland}{}{}{}{}
\cventry{March 2026}{Towards the Full Bloom of First Star, Galaxy,
and Black Hole Formation Exploration, Tokyo, Japan}{}{}{}{}
\cventry{April 2025}{Escape of Lyman radiation from galactic labyrinths, Crete, Greece}{}{}{}{}

\cventry{April 2024}{STScI Spring Symposium, Baltimore}{}{}{}{}


\cventry{Feb. 2024}{Building Galaxies From Scratch, Vienna}{}{}{}{}
\cventry{Nov. 2022}{CIERA Theory Seminar, Northwestern University}{}{}{}{}
\cventry{Nov. 2022}{CCAPP Seminar, Ohio State University}{}{}{}{}
\cventry{Oct. 2022}{Coffee Seminar, Institute for Advanced Study}{}{}{}{}
\cventry{Oct. 2022}{Cosmology Seminar, MPA Garching}{}{}{}{}
\cventry{Aug. 2022}{Star Formation in Different Environments, Quy Nhon, Vietnam}{}{}{}{}
\cventry{July 2022}{A Holistic View of Stellar Feedback and Galaxy Evolution, Ascona}{}{}{}{}
\cventry{July 2022}{Planet and Star Formation (PSF) Seminar - MPIA, Heidelberg}{}{}{}{}
\cventry{June 2022}{NBIA Workshop on Radiation Transfer in Astrophysics, Copenhagen}{}{}{}{}
\cventry{Nov. 2021}{IAU Symposium 362 - The Predictive Power of Computational Astrophysics}{}{}{}{}


\section{List of Publications}
% publications.tex
% Auto-generated by fetch_ads_pubs.py — do not edit by hand.
% Last synced: 2026-06-04

\textbf{\textit{33 publications: 10 first-author, 5 second-author; $h$-index: 15; total citations: 628
(\href{https://ui.adsabs.harvard.edu/public-libraries/2ISqHbaiQ9-l2n4w4stzSw}{ADS listing})}}\newline

\underline{\textbf{\textit{First-author publications}}} \newline
\begin{enumerate}

\item
    \href{https://doi.org/10.48550/arXiv.2610.17577}{{Non-equilibrium chemistry and radiation hydrodynamics in metal-poor and metal-rich interstellar gas}} \\
    \textbf{Menon, S.}, Sharda, P., Burkhart, B., Grassi, T. \textcolor{RedViolet}{{ApJ subm., 2026, arXiv:2610.17577}}
    
\item
    \href{https://arxiv.org/abs/2605.13982}{{Lyman-alpha Radiation Pressure in Dense Star Clusters: Implications for Star Formation and Winds at Cosmic Dawn}} \\
    \textbf{Menon, S.}, Smith, A. \textcolor{RedViolet}{{ApJ accepted (in press), 2026, arXiv:2605.13982}. } 

    \item
    \href{https://arxiv.org/abs/2408.14591}{{Bursts of Star Formation and Radiation-driven Outflows Produce Efficient LyC Leakage from Dense Compact Star Clusters}} \\
    \textbf{Menon, S.}, Burkhart, B., Somerville, R., Thompson, T., Sternberg, A. \textcolor{RedViolet}{{ApJ, 2025, 987, 12}}

    \item
    \href{https://arxiv.org/abs/2405.00813}{{The Interplay between the Initial Mass Function and Star Formation Efficiency through Radiative Feedback at High Stellar Surface Densities}} \\
    \textbf{Menon, S.}, Lancaster, L., Burkhart, B., Somerville, R., Dekel, A., Krumholz, M. \textcolor{RedViolet}{{ApJL, 2024, 967, L28}}

    \item
    \href{https://arxiv.org/abs/2210.02818}{{Outflows driven by direct and reprocessed radiation pressure in massive star clusters}} \\
    \textbf{Menon, S.}, Federrath, C., Krumholz, M. \textcolor{RedViolet}{{MNRAS, 2023, 521, 5160}}

    \item
    \href{https://arxiv.org/abs/2206.14190}{{Infrared radiation feedback does not regulate star cluster formation}} \\
    \textbf{Menon, S.}, Federrath, C., Krumholz, M. \textcolor{RedViolet}{{MNRAS, 2022, 517, 1313}}

    \item
    \href{https://arxiv.org/abs/2202.08778}{{VETTAM: a scheme for radiation hydrodynamics with adaptive mesh refinement using the variable Eddington tensor method}} \\
    \textbf{Menon, S.}, Federrath, C., Krumholz, M., Kuiper, R., Wibking, B., Jung, M. \textcolor{RedViolet}{{MNRAS, 2022, 512, 401}}

    \item
    \href{https://arxiv.org/abs/2108.04387}{{The dependence of the hierarchical distribution of star clusters on galactic environment}} \\
    \textbf{Menon, S.}, Grasha, K., Elmegreen, B., Federrath, C., Krumholz, M., Calzetti, D., et al. \textcolor{RedViolet}{{MNRAS, 2021, 507, 5542}}

    \item
    \href{https://arxiv.org/abs/2010.09861}{{On the compressive nature of turbulence driven by ionizing feedback in the pillars of the Carina Nebula}} \\
    \textbf{Menon, S.}, Federrath, C., Klaassen, P., Kuiper, R., Reiter, M. \textcolor{RedViolet}{{MNRAS, 2021, 500, 1721}}

    \item
    \href{https://arxiv.org/abs/2002.08707}{{On the turbulence driving mode of expanding H II regions}} \\
    \textbf{Menon, S.}, Federrath, C., Kuiper, R. \textcolor{RedViolet}{{MNRAS, 2020, 493, 4643}} \newline

\end{enumerate}

\underline{\textbf{\textit{Co-author publications with significant contributions}}} \newline

\begin{enumerate}

 \item
    \href{https://arxiv.org/abs/2604.00197}{{Radiative Feedback in Population III Protostellar Growth: The Impact of HI Shielding}} \\
    Chen, A., \textbf{Menon, S.}, Burkhart, B., Sharda, P., Williams, C., Naoz, S., et al. \textcolor{RedViolet}{{arXiv e-prints, 2026, arXiv:2604.00197}}

    \item
    \href{https://arxiv.org/abs/2503.11877}{{Planets and planetesimals at cosmic dawn: vortices as planetary nurseries}} \\
    Eriksson, L., \textbf{Menon, S.}, Carrera, D., Lyra, W., Burkhart, B. \textcolor{RedViolet}{{MNRAS, 2025, 542, 641}}

    \item
    \href{https://arxiv.org/abs/2501.12734}{{Magnetic fields limit the mass of Population III stars even before the onset of protostellar radiation feedback}} \\
    Sharda, P., \textbf{Menon, S.}, Gerasimov, R., Bromm, V., Burkhart, B., Haemmerlé, L., et al. \textcolor{RedViolet}{{MNRAS, 2025, 541, L1}}

    \item
    \href{https://arxiv.org/abs/2405.18265}{{Population III star formation in the presence of turbulence, magnetic fields, and ionizing radiation feedback}} \\
    Sharda, P., \textbf{Menon, S.} \textcolor{RedViolet}{{MNRAS, 2025, 540, 1745}}

     \item
    \href{https://arxiv.org/abs/2109.03983}{{First extragalactic measurement of the turbulence driving parameter: ALMA observations of the star-forming region N159E in the Large Magellanic Cloud}} \\
    Sharda, P., \textbf{Menon, S.}, Federrath, C., Krumholz, M., Beattie, J., Jameson, K., et al. \textcolor{RedViolet}{{MNRAS, 2022, 509, 2180}}

    \item
    \href{https://arxiv.org/abs/2602.18556}{{Blue Monsters and Dusty Descendants: Reconciling UV and IR Emission from Galaxies from z$\sim$7, up to z = 14}} \\
    Sommovigo, L., Lancaster, L., \textbf{Menon, S.}, O'Leary, J., Somerville, R., Bryan, G. \textcolor{RedViolet}{{ApJ, 2026, 1003, 170}}

    \item
    \href{https://arxiv.org/abs/2505.05442}{{Density-modulated star formation efficiency: implications for the observed abundance of ultraviolet luminous galaxies at $z>10$}} \\
    Somerville, R., Yung, L., Lancaster, L., \textbf{Menon, S.}, Sommovigo, L., Finkelstein, S. \textcolor{RedViolet}{{MNRAS, 2025, 544, 3774}}

    \item
    \href{https://arxiv.org/abs/2511.11314}{{Radiation magnetohydrodynamics simulations of Population III star formation during the epoch of reionization}} \\
    van Veenen, L., Sharda, P., Viti, S., \textbf{Menon, S.} \textcolor{RedViolet}{{A\&A, 2026, 708, A61}}
    
    \item
    \href{https://arxiv.org/abs/2505.22730}{{The Coevolution of Stellar Wind-blown Bubbles and Photoionized Gas. I. Physical Principles and a Semianalytic Model}} \\
    Lancaster, L., Kim, J., Bryan, G., \textbf{Menon, S.}, Ostriker, E., Kim, C. \textcolor{RedViolet}{{ApJ, 2025, 989, 42}} 
     \item
    \href{https://arxiv.org/abs/2609.01065}{{Investigating the Young Stellar Populations and Hierarchies in Nearby Galaxies with the UVIT. III. Evidence for a Largest Scale of Correlated Stellar Structures and a Non-universal Fractal Dimension}} \\
    Shashank, G., Subramanian, S., \textbf{Menon, S.}, Mondal, C., Grasha, K., de Grijs, R., et al. \textcolor{RedViolet}{{arXiv e-prints, 2026, arXiv:2609.01065}}

    \item
    \href{https://arxiv.org/abs/2502.00165}{{Linking Stellar Populations to H II Regions across Nearby Galaxies. II. Infrared Reprocessed and UV Direct Radiation Pressure in H II Regions}} \\
    Pathak, D., Leroy, A., Thompson, T., Lopez, L., Barnes, A., Dale, D., \textbf{incl. Menon, S.}, et al. \textcolor{RedViolet}{{ApJ, 2025, 982, 140}}

\end{enumerate}

\vspace{5pt}
\underline{\textbf{\textit{Contributions to collaborative efforts}}}
\vspace{5pt}
\begin{enumerate}
    \item\href{https://doi.org/10.48550/arXiv.2610.20719}{{Wandering intermediate-mass black holes and off-nuclear tidal disruption events from early star clusters}} \\
    Williams, C., Naoz, S., Zhou, Y., Akiba-Lewis, T., Vogelsberger, M., Burkhart, B., \textbf{incl. Menon, S.}, et al. \textcolor{RedViolet}{{arXiv e-prints, 2026, arXiv:2610.20719}}

    \item
    \href{https://doi.org/10.48550/arXiv.2610.30265}{{The THESAN-ZOOM project: the delayed rise of Lyman continuum escape fractions in bursty high-redshift galaxies}} \\
    Mehrotra, P., Smith, A., Bian, Y., McClymont, W., Puchwein, E., Kannan, R., \textbf{incl. Menon, S.}, et al. \textcolor{RedViolet}{{arXiv e-prints, 2026, arXiv:2610.30265}}

    \item
    \href{https://arxiv.org/abs/2606.12254}{{Investigating the young stellar populations and hierarchies in nearby galaxies with the UVIT. II. Presenting the properties of ~25,000 UV-detected star-forming clumps}} \\
    Shashank, G., Subramanian, S., Mondal, C., \textbf{Menon, S.}, Subramaniam, A. \textcolor{RedViolet}{{arXiv e-prints, 2026, arXiv:2606.12254}}
    \item
    \href{https://arxiv.org/abs/2412.00872}{{Tracing hierarchical star formation out to kiloparsec scales in nearby spiral galaxies with UVIT}} \\
    Shashank, G., Subramanian, S., Muraleedharan, S., \textbf{Menon, S.}, Mondal, C., Krishna, S., et al. \textcolor{RedViolet}{{A\&A, 2025, 693, A188}}
     \item
    \href{https://arxiv.org/abs/2602.22860}{{An FUV-optical Approach for Studying Hierarchical Star Formation in Nearby Galaxies with UVIT}} \\
    Ananthu, S., Shashank, G., Subramanian, S., Jayanth, R., \textbf{Menon, S.}, Mondal, C., et al. \textcolor{RedViolet}{{ApJ, 2026, 1002, 220}}
    \item
    \href{https://arxiv.org/abs/2603.26872}{{A cosmological framework for stellar collisions at high redshift in proto-globular clusters, nuclear star clusters, and Little Red Dots}} \\
    Williams, C., Naoz, S., Rose, S., Burkhart, B., Yoshida, N., Chen, A., \textbf{incl. Menon, S.}, et al. \textcolor{RedViolet}{{arXiv e-prints, 2026, arXiv:2603.26872}}
    \item
    \href{https://arxiv.org/abs/2509.22776}{{Observing Compact Population III Star Clusters and the Presence of Cosmic Streaming}} \\
    Williams, C., Naoz, S., Yoshida, N., Lake, W., Burkhart, B., Marinacci, F., \textbf{incl. Menon, S.}, et al. \textcolor{RedViolet}{{ApJ, 2026, 997, 202}}
    \item
    \href{https://arxiv.org/abs/2507.12993}{{Supersonic Turbulence in Primordial Halos: a Comparison with and without the Stream Velocity}} \\
    Chen, A., Lake, W., Williams, C., Burkhart, B., Naoz, S., \textbf{Menon, S.}, et al. \textcolor{RedViolet}{{ApJ, 2025, 995, 65}}
    \item
    \href{https://arxiv.org/abs/2502.19484}{{Detecting Molecular Hydrogen (H$_2$) Emission at Cosmic Dawn}} \\
    Johnson, M., Burkhart, B., D'Eugenio, F., Tacchella, S., Maiolino, R., Bialy, S., \textbf{incl. Menon, S.}, et al. \textcolor{RedViolet}{{ApJ, 2025, 992, 196}}
    \item
    \href{https://arxiv.org/abs/2502.17561}{{ΛCDM Star Clusters at Cosmic Dawn: Stellar Densities, Environment, and Equilibrium}} \\
    Williams, C., Naoz, S., Lake, W., Burkhart, B., Marinacci, F., Vogelsberger, M., \textbf{incl. Menon, S.}, et al. \textcolor{RedViolet}{{ApJ, 2025, 990, 135}}

    \item
    \href{https://arxiv.org/abs/2406.12025}{{A 260 pc resolution ALMA map of HCN(1$\rightarrow$0) in the galaxy NGC 4321}} \\
    Neumann, L., Bigiel, F., Barnes, A., Gallagher, M., Leroy, A., Usero, A., \textbf{incl. Menon, S.}, et al. \textcolor{RedViolet}{{A\&A, 2024, 691, A121}}

    \item
    \href{https://arxiv.org/abs/2401.14453}{{Hidden Gems on a Ring: Infant Massive Clusters and Their Formation Timeline Unveiled by ALMA, HST, and JWST in NGC 3351}} \\
    Sun, J., He, H., Batschkun, K., Levy, R., Emig, K., Rodríguez, M., \textbf{incl. Menon, S.}, et al. \textcolor{RedViolet}{{ApJ, 2024, 967, 133}}

\end{enumerate}
\vspace{5pt}
\underline{\textbf{\textit{Popular science articles}}} \newline

\begin{enumerate}
    \item \href{https://physicstoday.aip.org/features/how-stars-shape-galaxies}{{How stars shape galaxies}} \\
    Orr, M., \textbf{Menon, S.}, Burkhart, B. \textcolor{RedViolet}{{Physics Today, 2025, Issue 5, id. 44279}}
\end{enumerate}



% \section{Selected Talks}
% 2 invited talks, 15 contributed talks
% \cventry{October 2022}{CCAPP Seminar, Ohio State University}{\newline Title: Role of stellar feedback in extreme environments}{}{}{}
% \cventry{October 2022}{Cosmology Seminar, MPA Garching}{\newline Title: Stellar feedback in the formation of dense stellar systems}{}{}{}
% \cventry{September 2022}{SFIR Seminar, Princeton University}{\newline Title: Radiation Feedback in Star Cluster Formation and Evolution}{}{}{}
% \cventry{August 2022}{Contributed Talk - Star Formation in Different Environments, Quy Nhon, Vietnam}{\newline Title: The dynamical role of radiation pressure in the turbulent ISM}{}{}{}
% \cventry{July 2022}{Contributed Talk - A Holistic View of Stellar Feedback and Galaxy Evolution, Ascona}{\newline Title: The role of radiation pressure in dense stellar systems}{}{}{}
% \cventry{July 2022}{Planet and Star Formation (PSF) Seminar - Max-Planck-Institut for Astronomy, Heidelberg}{\newline Title: Radiation Feedback in Star Cluster Formation}{}{}{}
% \cventry{June 2022}{Seminar - Centre de Recherche Astrophysique de Lyon (CRAL)}{\newline Title: Radiation Feedback in Star Cluster Formation}{}{}{}
% \cventry{June 2022}{Star Formation Seminar - CEA, Saclay}{\newline Title: Radiation Feedback in Star Cluster Formation}{}{}{}
% \cventry{June 2022}{Poster - From Stars to Galaxies II, Gothenburg}{\newline Title: The inability of radiation pressure to regulate star formation}{}{}{}
% \cventry{June 2022}{Contributed Talk - NBIA Workshop on Radiation Transfer in Astrophysics, Copenhagen}{\newline Title: Modelling radiation hydrodynamics with the VET method}{}{}{}
% \cventry{Dec 2021}{Invited Colloquium - Indian Institute of Astrophysics, Bangalore}{\newline Title: The interplay of stellar feedback and turbulence in the ISM}{}{}{}
% \cventry{Nov 2021}{Contributed Talk, IAU Symposium 362 - The Predictive Power of Computational Astrophysics}{\newline Title: \texttt{VETTAM} - A novel algorithm for modelling radiation hydrodynamics}{}{}{}
% \cventry{May 2021}{Contributed Talk, ISM 2021, Beirut - Structure, Characteristic Scales, and Star Formation}{\newline Title: Spatial distribution of star clusters in external galaxies}{}{}{}
% \cventry{February 2021}{Contributed Talk, Australian National Insitute for Theoretical Astrophysics science workshop}{\newline Title: Compressive Turbulence Driven by Ionising Radiation in the ISM}{}{}{}
% \cventry{September 2020}{Contributed Talk, Annual meeting of the Astronomische Gesellschaft, Germany}{\newline Title: Compressive Turbulence Driven by Ionising Radiation}{}{}{}
% \cventry{March 2020}{Invited Talk, Modelling High-Mass Stellar Feedback Workshop, T\"{u}bingen, Germany}{\newline Title: Turbulence and Triggered Star Formation in Star Clusters}{}{}{}
% \cventry{January 2020}{Australian National Insitute for Theoretical Astrophysics science workshop}{\newline Title: Turbulence Driving Mode of Expanding HII regions}{}{}{}

\end{document}
